% !PRINTED 62
\weberpage{62}

The function $\Phi(x,\xi)$ is called the \gesperrt{polar} of $\Phi$. It is linear and homogeneous both with respect to the $x$ and with respect to the $\xi$.

It can be expressed by
\begin{equation}
 \Phi(x,\xi)=2\sum a_{ik}\xi_i x_k
\tag{14}
\end{equation}
and satisfies the condition
\begin{equation}
 \Phi(x,x)=2\Phi(x).
\tag{15}
\end{equation}

\begin{center}
{\large \S.\ 17.}

\vspace{0.5em}
{\large\gesperrt{Euler's Theorem on Homogeneous Functions.}}
\end{center}

From the preceding developments one can easily derive a fundamental theorem on homogeneous functions, discovered by Euler and named after him.

We obtain it most simply from formula (4) of the preceding paragraph if, denoting by $t$ an arbitrary variable, we put
\begin{equation}
 \xi_1=tx_1,\quad \xi_2=tx_2,\ldots,\quad \xi_m=tx_m
\tag{1}
\end{equation}
and then apply the fundamental formula \S~15 (1) for homogeneous functions. We thus obtain first:
\begin{equation}
 (1+t)^n\Phi(x)
 =
 \sum { (tx_1)^{\alpha_1}(tx_2)^{\alpha_2}\cdots(tx_m)^{\alpha_m}
 \over \Pi(\alpha_1)\Pi(\alpha_2)\cdots\Pi(\alpha_m)}
 D_{\alpha_1,\alpha_2,\ldots,\alpha_m}\Phi.
\tag{2}
\end{equation}
If one applies the binomial theorem to the left-hand side of (2) and then equates on both sides the coefficients of equal powers of $t$, one obtains for every $\nu=1,2,\ldots,n$:
\begin{equation}
 {\Pi(n)\over\Pi(n-\nu)}\Phi(x_1,x_2,\ldots,x_m)
 =
 \sum^{\alpha}
 {\Pi(\nu)\over\Pi(\alpha_1)\Pi(\alpha_2)\cdots\Pi(\alpha_m)}
 x_1^{\alpha_1}x_2^{\alpha_2}\cdots x_m^{\alpha_m}
 D_{\alpha_1,\alpha_2,\ldots,\alpha_m}\Phi,
\tag{3}
\end{equation}
where the sum is extended over all systems of values of the $\alpha$ satisfying the condition
\begin{equation}
 \alpha_1+\alpha_2+\cdots+\alpha_m=\nu.
\tag{4}
\end{equation}

In this form the theorem to be proved is contained in its full generality. For the special case $\nu=1$ we obtain the formula
%CONT>>
